\scanpage{208}
\mspage{157}

\[
\alpha_1(x_0,u)\in W\quad\text{\unsure{i.e.}}\quad (X_{0/x_0})_u\quad W\text{-asphérique}.
\]

\deleted{Cela signifie} Dire que c'est vrai pour tt $u$ dans \deleted{\unsure{la}} $Y_{/y}$, signifie que $f_0\colon X_0\unsure{-|\longrightarrow}Y$ est lisse en \unsure{$x$|$x_0$|$u$}.

Donc on a prouvé a), sous la forme plus précise :

\uline{Corollaire 6} \deleted{\unsure{en supposant $W(7\ \mathrm{bis})$|On suppose $W(7\ \mathrm{bis})$}} Soit $x\in\operatorname{Ob}\mathcal X$ (où

\diagram{\begin{tikzcd}[ampersand replacement=\&,row sep=0pt,column sep=small]
\& X_1\\
\mathcal X=\displaystyle\int X_0 \arrow[ur,dash]\arrow[dr] \&\\
\& X_2
\end{tikzcd}}

). On veut exprimer la \inserted{$W$-}lissité de \deleted{\unsure{$\varphi$}}

\[
\varphi_Q\colon\mathcal X\unsure{-|\longrightarrow}Y
\]

en $x$.

a) Si $x\in\operatorname{Ob}X_i$ $(i\in\{1,2\})$, $\varphi_Q$ est \inserted{$W$-}lisse en $x$ ssi $f_i\colon X_i\unsure{-|\longrightarrow}Y$ l'est.

b) Supposons $W(7\ \mathrm{bis})$\textsuperscript{*}. \deleted{Si $x\in\operatorname{Ob}X_0$,} \inserted{Soit $x_0\in\operatorname{Ob}X_0$, et supposons $f_1,f_2$ lisses en $x_1,x_2$.} Alors $\varphi_Q$ est $W$-lisse en \unsure{$x_0$} ssi $f_0\colon X_0\longrightarrow Y$ l'est.

% slanted
\margin{Inutile pour la partie directe, en supposant}

Donc \inserted{(moyennant $W(7\ \mathrm{bis})$)} dans \deleted{les} les trois cas $x\in\operatorname{Ob}X_i$ $(i\in\{0,1,2\})$, on trouve que $\varphi_Q$ est \inserted{$W$-}lisse en $x$ ssi $f_i\colon X_i\unsure{-|\longrightarrow}Y$ l'est.

Quant à b) du \deleted{\unsure{la prop. II|le princ. II|la prop. 15}}, \inserted{th. \deleted{\illegible{1}}} c'est le lemme 2 ci-dessus. Donc tout est prouvé !
% archive sheet number: 208
